The previous chapters produce, for every small diffusivity $\eps$, an eigenpair
$(\lambda^\eps,h^\eps)$ of the \emph{horizontal} operator $A_{\eps,K}$ on
$\Lp^2(\T^2;\CC^2)$, with $\re\lambda^\eps\geq g_*>0$.  Two components are not a magnetic
field.  This chapter recovers the third one: the divergence constraint determines it
algebraically, $b_3=-\Div_\perp h/(2\pi iK)$, and its own equation supplies the regularity
needed for membership in the operator domain (Lemma~\ref{lem:axial-recovery}).  The
resulting $V^\eps=\e^{2\pi iK_\eps x_3}b^\eps$ is divergence-free \emph{by construction},
which is why the paper's proof never needs a separate solenoidality lemma.  The remaining
two sections are independent of that: the non-zero axial Fourier index lets one
pass to real parts without losing any norm, at any temporal phase, so that the real datum
$B_{\rm in}^\eps=\re V^\eps$ is non-zero and its evolution obeys the exact exponential law;
and a last section records the bounds behind the two-sided $\sqrt\eps$ magnetic scale.

\textbf{Standing conventions of the formalisation.}  Three global conventions apply to every node
of this chapter and are not repeated in each one.  (i) The three-torus is modelled as the product
$\T^2\times\T^1$ (\texttt{ProductTorus3}) with the product measure; a measure-preserving
equivalence with the literal $\T^3=(\R/\Z)^3$ and the induced isometric $\Lp^2$ transport do exist
(\texttt{FastDynamo.productToLiteral}, \texttt{FastDynamo.toLiteralR3L2} in
\texttt{FastDynamo/Analysis/LiteralThreeTorus.lean}), and the real magnetic trajectory and its norm
law are transported along them, but no differential operator is transported, so all operator
statements below live on the product model.  (ii) The velocity is an element of
\texttt{TransverseRealW1Inf}, the Fourier-graph realisation of $\W^{1,\infty}(\T^2;\R^3)$; such a
field is independent of $x_3$ but has a non-trivial third component.  This is exactly the class of
the velocity constructed in the previous chapters, and it is what makes the axial reduction exact.
Axial independence is \emph{not} a narrowing: the only lemma of this section that hypothesises a
velocity at all, Lemma~\ref{lem:axial-recovery}, asks for a divergence-free
$u\in\W^{1,\infty}(\T^3;\R^3)$ \emph{independent of $x_3$} (as in the paper's statement of
Lemma~8.1), which is exactly this class; Lemma~\ref{lem:real-axial-mode} has no velocity, and
Corollary~\ref{cor:magnetic-scale} is about the theorem's own $u$.  What does remain
unformalised is the identification of the Fourier-graph model \texttt{TransverseRealW1Inf} with
the Lipschitz / $\W^{1,\infty}(\T^3)$ class itself (see the scope notes in
\texttt{README.md}).  Divergence-freeness of $u$
is \texttt{IsTransverseDivergenceFree}, the identity $\partial_1u_1+\partial_2u_2=0$ in
$\Lp^\infty$.  (iii) $\Hs^2$ and $\Hs^1$ domains are Fourier graph spaces, and all weak
derivatives, divergences and curls are their Fourier realisations.

\section{Recovery of the axial magnetic component}

\begin{lemma}[Recovery of the axial magnetic component]
  \label{lem:axial-recovery}
  \lean{FastDynamo.axialRecoveredComponent,
    FastDynamo.horizontalEmbedH2, FastDynamo.transverseSnocH2,
    FastDynamo.transverseCoordL2_reducedInductionAction,
    FastDynamo.fixedSectorDivergenceCoeff_reducedInductionAction,
    FastDynamo.axialRecoveredSymbol,
    FastDynamo.transverseH2Weight_le_axialRecoveredSymbol,
    FastDynamo.axialRecoveredCertificateMap,
    FastDynamo.axialRecoveredComponent_coeff_equation,
    FastDynamo.axialRecoveredCertificate_eq,
    FastDynamo.exists_reducedInductionEigenpair_of_horizontalEigenpair}
  \leanok
  \uses{def:operator:horizontal-induction, def:main:induction-operator}
  Let $\eps>0$, $K\in\Z\setminus\{0\}$, $\lambda\in\CC$, and let
  $u\in\W^{1,\infty}(\T^3;\R^3)$ be divergence-free and independent of $x_3$.  Suppose
  $0\neq h\in\Hs^2(\T^2;\CC^2)$ satisfies $A_{\eps,K}h=\lambda h$.  Put
  \[
    b_3=-\frac{\Div_\perp h}{2\pi iK},\qquad b=(h,b_3),
    \qquad V(y,x_3)=\e^{2\pi iKx_3}b(y).
  \]
  Then $b\in\Hs^2(\T^2;\CC^3)$ and
  \[
    0\neq V\in\Hs^2(\T^3;\CC^3)\cap\Lp^2_\sigma(\T^3),\qquad L_{\eps,u}V=\lambda V .
  \]
\end{lemma}

Three points of the Lean form.  (a) The conclusion is packaged as an eigenpair of the
\emph{three}-component reduced operator $A^{(3)}_{\eps,K}$ of
Definition~\ref{def:main:induction-operator} together with
$\pi_\perp b=h$, the formula for the third coordinate, and \emph{both} divergence clauses
$\Div_\perp h+2\pi iKb_3=0$ and $\Div V=0$; the passage to the physical curl form is then
Lemma~\ref{lem:finish:physical-compatibility}.  (b) The identity
\eqref{eq:finish:transport-stretching-divergence} is not re-proved in two dimensions: the
product-torus identity
$\Div[-(w\cdot\nabla)F+(F\cdot\nabla)w]=-w\cdot\nabla\Div F$ is transported into the sector
$K$ along the axial lift, which avoids a second derivative of $w$.  The paper does not
use that identity (see the deviation below).  (c) The elliptic upgrade $b_3\in\Hs^1\Rightarrow b_3\in\Hs^2$ is
\emph{free} at every frequency, with no elliptic-regularity theorem: the multiplier
$(1+4\pi^2\abs n^2)/\bigl(\eps(4\pi^2\abs n^2+(2\pi K)^2)\bigr)$ has modulus at most
$1/\eps$ as soon as $(2\pi K)^2\geq1$, which holds for every $K\neq0$, so the $\Hs^2$
certificate is a bounded Fourier multiplier.

\textbf{Deviation: the scalar route to $b_3$, not the three-component distributional
route.}  The statement above is the paper's (Lemma~8.1, \texttt{lem:axial-recovery}), word for
word.  The proof below is not.  The paper proves the lemma without writing an equation for
$b_3$: for the $\Hs^1$ field $V$, which is divergence-free by construction, it forms
$F=\lambda V-\eps\Delta V-\curl(u\times V)\in\Hs^{-1}(\T^3;\CC^3)$, uses the identity
$\curl(u\times V)=(V\cdot\nabla)u-(u\cdot\nabla)V$ in distributions, observes that $F_\perp=0$ is
the horizontal eigenvalue equation and that $\Div F=0$, so that, every term having axial
Fourier index $K$, $0=\Div F=\partial_{x_3}F_3=2\pi iKF_3$ in distributions forces $F=0$; only
then does it upgrade $V$ to $\Hs^2$ by elliptic regularity from
$\eps\Delta V=\lambda V-\curl(u\times V)\in\Lp^2$.  Lean takes a scalar route instead: the
transport--stretching divergence identity
$\Div[-(w\cdot\nabla)F+(F\cdot\nabla)w]=-w\cdot\nabla\Div F$, transported into the sector
$K$ along the axial lift, turns $A_{\eps,K}h=\lambda h$ into the scalar equation
$(\lambda-S)b_3=(h\cdot\nabla_\perp)u_3$, whose right-hand side is an $\Lp^2$ field;
the $\Hs^2$ certificate for $b_3$ is then a bounded Fourier multiple of it, with no
elliptic-regularity theorem.  The paper's route is not available as stated because Lean's
curl form is realised in $\Hs^{-2}$ against $\Hs^2$ test fields, not in $\Hs^{-1}$ against
$\Hs^1$: forming $F$ at the regularity the paper uses would require a further antidual layer
and the distributional identity $\curl(u\times V)=(V\cdot\nabla)u-(u\cdot\nabla)V$ one
derivative lower.  The two cheap steps of the paper's route -- $\Div F=0\Rightarrow
2\pi iKF_3=0\Rightarrow F=0$ in a single non-zero axial mode, and $\eps\Delta V\in\Lp^2
\Rightarrow V\in\Hs^2$ -- are free in the Fourier model, and are how point (c) above is
proved.  The statement proved is the paper's, including both divergence clauses; only the
route differs.  The paper has no counterpart of the two displays
\eqref{eq:finish:transport-stretching-divergence} and
\eqref{eq:finish:recovered-axial-equation}; their labels are blueprint-local.

\begin{proof}
  \leanok
  \uses{def:operator:horizontal-induction, def:main:induction-operator}
  This is the argument Lean states, not the paper's; see the deviation
  above.  Write $d=\Div_\perp h$ and
  $S=\eps\bigl(\Delta_\perp-(2\pi K)^2\bigr)-u_\perp\cdot\nabla_\perp-2\pi iKu_3$.  For a divergence-free
  $w\in\W^{1,\infty}$ and $F\in\Hs^2$,
  \begin{equation}
    \Div\bigl[-(w\cdot\nabla)F+(F\cdot\nabla)w\bigr]=-w\cdot\nabla\Div F
    \label{eq:finish:transport-stretching-divergence}
  \end{equation}
  in distributions, because the bracket is $\curl(w\times F)-w\Div F$; this needs no second
  derivative of $w$.  Applying it to the horizontal fields and using
  $\Div_\perp(u_3h)=u_3d+h\cdot\nabla_\perp u_3$ turns $A_{\eps,K}h=\lambda h$ into
  $(\lambda-S)d=-2\pi iK(h\cdot\nabla_\perp)u_3$ in $\Hs^{-1}(\T^2)$, that is
  \begin{equation}
    (\lambda-S)b_3=(h\cdot\nabla_\perp)u_3 .
    \label{eq:finish:recovered-axial-equation}
  \end{equation}
  Initially $b_3\in\Hs^1$; rearranging \eqref{eq:finish:recovered-axial-equation} gives
  $\eps\Delta_\perp b_3\in\Lp^2$ and hence $b_3\in\Hs^2$.  The horizontal equation together
  with \eqref{eq:finish:recovered-axial-equation} is exactly
  $\eps\Delta V-(u\cdot\nabla)V+(V\cdot\nabla)u=\lambda V$, and
  $\Div V=\e^{2\pi iKx_3}(d+2\pi iKb_3)=0$, so the left-hand side is $L_{\eps,u}V$.  Finally $h\neq0$
  gives $V\neq0$.
\end{proof}

\section{The eigenfunction is divergence-free}

For the recovered mode, $\Div V=0$ holds by construction (Lemma~\ref{lem:axial-recovery}).
The lemma of this section proves it for every eigenfunction of the auxiliary operator
$\mathcal{A}_\eps$ with $\re\lambda>0$; the Lean proof of
Proposition~\ref{prop:finish:real-strong-mode} takes the solenoidality of the lifted mode
from it, and the dynamics chapter uses its resolvent form
(Lemma~\ref{lem:dynamics:solenoidal-subspace}).  The passage from $\mathcal{A}_\eps$ to the
physical curl form (Lemma~\ref{lem:finish:physical-compatibility}) is used by the axial
recovery and by the dynamics chapter.

\begin{definition}[The physical curl-form operator]
  \label{def:finish:physical-operator}
  \lean{FastDynamo.productL2ToHMinus2, FastDynamo.productWeakCurl,
  FastDynamo.productCrossWithVelocity, FastDynamo.productPhysicalInductionAction,
  FastDynamo.productPhysicalInductionOperator}
  \leanok
  \uses{def:main:induction-operator}
  For $u\in\W^{1,\infty}$ and $\eps\in\R$ the \emph{physical} induction operator is the curl form
  \[
    L_{\eps,u}B=\eps\Delta B+\curl(u\times B),
  \]
  realised as a bounded map from the Fourier graph $\Hs^2(\T^2\times\T^1;\CC^3)$ into a Fourier
  model of $\Hs^{-2}$, in which an $\Lp^2$ coefficient $\hat f(n,k)$ is represented by
  $\hat f(n,k)/(1+\abs\xi^2)$.  In that model the distributional curl of an arbitrary $\Lp^2$
  field is bounded, and $\Div\circ\curl=0$ is an exact identity frequency by frequency.
\end{definition}

In Lean this is \texttt{productPhysicalInductionAction} and the bundled domain operator
\texttt{productPhysicalInductionOperator}, whose ambient space is $\Hs^{-2}$ and whose embedding
is $\Hs^2\hookrightarrow\Lp^2\hookrightarrow\Hs^{-2}$ (module
\texttt{FastDynamo/Operator/Compatibility.lean}).  The paper never names a functional-analytic
realisation of $L_{\eps,u}$; the weak $\Hs^{-2}$ realisation is the formalisation's choice, and
it is weaker than an $\Lp^2$ realisation on solenoidal fields, which does not exist in the
development ($\Lp^2_\sigma$ is a predicate, not a subspace).

\begin{lemma}[Compatibility of the two actions]
  \label{lem:finish:physical-compatibility}
  \lean{FastDynamo.productInductionAction_toHMinus2_eq_physical,
  FastDynamo.productInductionEigenpair_to_physical}
  \leanok
  \uses{def:main:induction-operator, def:finish:physical-operator}
  Let $u$ be divergence-free and let $V\in\Hs^2$ satisfy $\Div V=0$.  Then
  \[
    \curl(u\times V)=-(u\cdot\nabla)V+(V\cdot\nabla)u,
  \]
  so $\mathcal{A}_\eps V$ and $L_{\eps,u}V$ agree in $\Hs^{-2}$.  Consequently an eigenpair of
  $\mathcal{A}_\eps$ whose eigenfunction is divergence-free is an eigenpair of $L_{\eps,u}$.
\end{lemma}

Faithful, modulo the standing conventions.

\begin{proof}
  \leanok
  \uses{def:main:induction-operator}
  The general vector identity for divergence-free $u$ reads
  $-(u\cdot\nabla)V+(V\cdot\nabla)u=\curl(u\times V)-u\,\Div V$, and $\Div V=0$ kills the last
  term.  In Lean the identity is proved coefficientwise from weak Leibniz certificates
  (\texttt{productInductionLowerField\_add\_velocityDiv\_eq\_weakCurl}), and the eigenpair transfer
  is then immediate because the embedding $\Lp^2\to\Hs^{-2}$ is injective.
\end{proof}

\begin{lemma}[The divergence equation and its energy]
  \label{lem:finish:divergence-energy}
  \lean{FastDynamo.productInductionLowerField_add_velocityDiv_eq_weakCurl,
  FastDynamo.productWeakDivergence_productInductionAction,
  FastDynamo.productWeakDivergence_velocityTimesScalar_pair_re_eq_zero,
  FastDynamo.productScalarWeakLaplacian_pair_self_re_nonpos,
  FastDynamo.productScalarWeakEigenfunction_eq_zero}
  \leanok
  \uses{def:main:induction-operator}
  Let $u\in\W^{1,\infty}$ be divergence-free, let $V\in\Hs^2$ and put $d=\Div V\in\Hs^1$.  Then, in
  the sense of distributions,
  \[
    -(u\cdot\nabla)V+(V\cdot\nabla)u+u\,d=\curl(u\times V),
    \qquad
    \Div\left[\mathcal{A}_\eps V\right]=\eps\Delta d-\Div(u\,d).
  \]
  Moreover the transport term is skew and the Laplacian is dissipative,
  \[
    \re\inner{\Div(u\,d)}{d}=0,
    \qquad
    \re\inner{\Delta d}{d}=-\norm{\nabla d}_2^2\leq0,
  \]
  so that a scalar $d\in\Hs^1$ obeying $\lambda d=-\Div(u\,d)+\eps\Delta d$ weakly, with
  $\eps\geq0$ and $\re\lambda>0$, vanishes.
\end{lemma}

In Lean the two pairings are the $\Hs^{-1}$--$\Hs^1$ duality, $\Div$ of an $\Lp^2$ field is
\texttt{productWeakDivergence} valued in $\Hs^{-1}$, and the final implication is the
velocity-free algebraic endpoint \texttt{productScalarWeakEigenfunction\_eq\_zero} (module
\texttt{FastDynamo/Analysis/ProductScalarEnergy.lean}); the vector identity and the divergence of
the induction action are in \texttt{FastDynamo/Analysis/InductionDivergence.lean}, the skewness in
\texttt{FastDynamo/Analysis/ProductTransportEnergy.lean}.

\begin{proof}
  \leanok
  The vector identity is checked on each of the three coordinates from the weak Leibniz rule for an
  $\Lp^\infty$ factor with an $\Lp^2$ factor; taking weak divergences and using
  $\Div\circ\curl=0$ gives the second display, since the diffusion term contributes
  $\eps\Delta d$.  Testing the scalar equation with $\bar d$, using periodicity and $\Div u=0$ for
  the transport term, and taking real parts yields
  $(\re\lambda)\norm d_2^2=-\eps\norm{\nabla d}_2^2\leq0$, whence $d=0$.
\end{proof}

\begin{lemma}[Eigenfunctions with $\re\lambda>0$ are divergence-free]
  \label{lem:positive-eigenfunction-solenoidal}
  \lean{FastDynamo.productInductionEigenpair_divergence_eq_zero,
  FastDynamo.productInductionEigenpair_solenoidal,
  FastDynamo.axiallyIndependent_positive_eigenfunction_solenoidal,
  FastDynamo.axiallyIndependent_eigenpair_solenoidal}
  \leanok
  \uses{def:main:induction-operator, def:finish:physical-operator}
  Let $u\in\W^{1,\infty}(\T^3;\R^3)$ be divergence-free and let $\eps\geq0$.  If
  $0\neq V\in\Hs^2(\T^3;\CC^3)$ satisfies
  \[
    \mathcal{A}_\eps V=\lambda V,
    \qquad\re\lambda>0,
  \]
  then $\Div V=0$.  Consequently $L_{\eps,u}V=\lambda V$.
\end{lemma}

The paper states no such lemma, because its recovered mode is solenoidal by construction
(the definition of $b_3$ in Lemma~8.1 of the paper); this is a result of the
formalisation.  The sign hypothesis on $\eps$ cannot be dropped: the statement is false for
$\eps<0$ (Lemma \ref{lem:finish:negative-diffusivity}).  The Lean statement therefore assumes
$0\leq\eps$, which is weaker than the paper's standing $\eps>0$ and hence covers it; the
positive-diffusivity wrapper is
\texttt{axiallyIndependent\_positive\_eigenfunction\_solenoidal}, and
\texttt{axiallyIndependent\_eigenpair\_solenoidal} also returns the physical eigenpair.  As explained in the
chapter preamble, the Lean velocity is $x_3$-independent; the general $\T^3$-velocity form of the
lemma is not formalised.  The conclusion comes in two equivalent shapes, $\Div V=0$ as an
$\Lp^2$ field and \texttt{IsProductSolenoidal} coefficientwise.

\begin{proof}
  \leanok
  \uses{lem:finish:divergence-energy, lem:finish:physical-compatibility}
  Set $d=\Div V$.  Taking the weak divergence of the eigenvalue equation and using the two displays
  of Lemma \ref{lem:finish:divergence-energy} gives the scalar advection-diffusion equation
  $\lambda d=-\Div(u\,d)+\eps\Delta d$ in $\Hs^{-1}$; since $d\in\Hs^1$ it is an admissible test
  function, and the energy endpoint of the same lemma gives $d=0$.  The last sentence is Lemma
  \ref{lem:finish:physical-compatibility}.  The argument starts from the auxiliary operator, not
  from the curl form whose agreement is only known after solenoidality, so the chain is not
  circular.
\end{proof}

\begin{lemma}[The sign hypothesis is necessary]
  \label{lem:finish:negative-diffusivity}
  \lean{FastDynamo.negativeDiffusivityEigenvalue,
  FastDynamo.negativeDiffusivity_auxiliary_eigenpair,
  FastDynamo.negativeDiffusivityEigenvalue_pos,
  FastDynamo.negativeDiffusivity_mode_not_productSolenoidal,
  FastDynamo.exists_positive_nonsolenoidal_eigenpair_of_negative_diffusivity}
  \leanok
  \uses{def:main:induction-operator}
  For every $\eps<0$ the zero velocity admits an eigenpair of $\mathcal{A}_\eps$ with positive real
  part and non-solenoidal eigenfunction: with $V(y,x_3)=\e^{2\pi ix_3}e_3$ one has
  \[
    \mathcal{A}_\eps V=\lambda V,
    \qquad\lambda=-\eps(2\pi)^2>0,
    \qquad\Div V\neq0 .
  \]
  Hence the hypothesis $\eps\geq0$ cannot be dropped from Lemma
  \ref{lem:positive-eigenfunction-solenoidal}.
\end{lemma}

This node is not in the paper; it is a remark-level Lean result recording that the sign
hypothesis of Lemma~\ref{lem:positive-eigenfunction-solenoidal} is sharp (module \texttt{FastDynamo/Finish/NegativeDiffusivity.lean}).  The main theorem is
unaffected, since it assumes $0<\eps$ throughout.

\begin{proof}
  \leanok
  With $u=0$ the transport and stretching terms vanish, and on the axial sector $K=1$ the constant
  vertical amplitude is annihilated by the transverse Laplacian, so only the axial diffusion symbol
  $-\eps(2\pi)^2$ survives.  Its divergence has the non-zero Fourier coefficient
  $2\pi i\cdot1$ at the zero transverse frequency.
\end{proof}

\section{Passage to a real-valued solution}

\begin{lemma}[Real part of a non-zero $x_3$-Fourier mode]
  \label{lem:real-axial-mode}
  \lean{FastDynamo.realPart_axialLiftPointwise_norm_sq,
  FastDynamo.phaseAxialFieldL2_norm_sq,
  FastDynamo.phaseAxialFieldL2_norm_sq_eq,
  FastDynamo.phaseAxialFieldL2_ne_zero}
  \leanok
  Let $K\in\Z\setminus\{0\}$, $b\in\Lp^2(\T^2;\CC^3)$ and $\varphi\in\R$.  Then, for almost every
  $y$,
  \[
    \int_0^1\abs{\re\bigl(\e^{i\varphi}\e^{2\pi iKx_3}b(y)\bigr)}^2\dd x_3
    =\frac12\abs{b(y)}_{\CC^3}^2 .
  \]
  In particular $\norm{\re(\e^{i\varphi}\e^{2\pi iKx_3}b)}_2^2=\frac12\norm b_2^2$ for every phase,
  this quantity does not depend on $\varphi$, and if $b\neq0$ the real part of the corresponding
  $x_3$-Fourier mode is non-zero for every phase.
\end{lemma}

Faithful.  The Lean statement comes in two layers: the fibrewise identity
\texttt{realPart\_axialLiftPointwise\_norm\_sq} is the displayed equation verbatim, for a
pointwise representative and the Haar measure of the axial circle; the global form
\texttt{phaseAxialFieldL2\_norm\_sq} is stated for the $\Lp^2$-class field
\texttt{phaseAxialFieldL2}~$K\,\varphi\,b=\re(\e^{i\varphi}\e^{2\pi iKx_3}b)$, which lives in the
real space $\Lp^2(\T^2\times\T^1;\R^3)$.  Module
\texttt{FastDynamo/Fourier/AxialL2.lean}.

\begin{proof}
  \leanok
  Write $b=a+ic$ with $a,c\in\R^3$.  Componentwise the integrand is the square of
  $a\cos(2\pi Kx_3+\varphi)-c\sin(2\pi Kx_3+\varphi)$.  Because $K$ is a non-zero integer, the
  squared sine and cosine have average $1/2$ on $[0,1]$ while their product has average zero;
  summing over the components gives the identity.  Integrating in $y$ by Fubini gives the global
  form.
\end{proof}

\begin{definition}[The real datum, the real solution and the strong-mode certificate]
  \label{def:finish:real-solution}
  \lean{FastDynamo.realAxialInitialField, FastDynamo.realAxialTrajectory,
  FastDynamo.productConjH2, FastDynamo.productRealPartH2,
  FastDynamo.complexifyProductR3L2, FastDynamo.realifiedEigenTrajectoryH2,
  FastDynamo.DomainOperator.IsStrongSolution,
  FastDynamo.ConditionalModeCertificate, FastDynamo.RealStrongModeCertificate}
  \leanok
  \uses{def:main:induction-operator, def:finish:physical-operator}
  Let $K\in\Z\setminus\{0\}$, $\lambda\in\CC$ and let $b$ be a transverse amplitude, and write
  $V=\e^{2\pi iKx_3}b$, $\alpha=\re\lambda$, $\omega=\im\lambda$.  Define
  \[
    B_{\rm in}=\re V,
    \qquad
    B(t)=\e^{\alpha t}\,\re\bigl(\e^{i\omega t}V\bigr)
    =\re\bigl(\e^{t\lambda}V\bigr)
    \qquad(t\in\R),
  \]
  both elements of the real space $\Lp^2(\T^2\times\T^1;\R^3)$
  (\texttt{realAxialInitialField}, \texttt{realAxialTrajectory}).  Inside the complex $\Hs^2$ graph
  the same field is represented by $W(t)=\frac12(\id+\mathrm{conj})(\e^{t\lambda}V)$
  (\texttt{realifiedEigenTrajectoryH2}, built from \texttt{productConjH2} and
  \texttt{productRealPartH2}), a fixed point of conjugation whose complexification is $B(t)$.  A
  curve $t\mapsto W(t)$ in the domain of a bundled operator $A$ is a \emph{strong solution} when
  $s\mapsto A.\text{embedding}(W(s))$ is differentiable in the ambient space with derivative
  $A.\text{op}(W(t))$ at every $t$ (\texttt{DomainOperator.IsStrongSolution}).  Finally
  \texttt{RealStrongModeCertificate}~$u\,\eps\,K\,\lambda\,b\,g$ bundles all the conclusions of
  this chapter for one mode: the lifted auxiliary eigenpair and its solenoidality, the physical
  eigenpair, $B_{\rm in}\neq0$, $B(t)\neq0$ for all $t$, the exact norm law, the lower bound with
  rate $g$, the conjugation-fixedness and solenoidality of $W(t)$, the identification of $W(t)$
  with $B(t)$, and the strong-solution equation for $W$.  (\texttt{ConditionalModeCertificate} is
  the sub-bundle not mentioning the real evolution equation.)
\end{definition}

The paper speaks of \emph{the} solution with datum $B_{\rm in}^\eps$, and its proof of
Theorem~1.2 writes that solution down.  This chapter exhibits the explicit strong solution in
$\Hs^2$ of the induction equation in the $\Hs^{-2}$ realisation of
Definition~\ref{def:finish:physical-operator}, with $W(0)=B_{\rm in}$; that it is the only one,
and that it is the orbit of $B_{\rm in}$ under the semigroup $\e^{tL_{\eps,u}}$ on
$\Lp^2_\sigma(\T^3)$, is Lemma~\ref{lem:dynamics:physical-uniqueness} in the appendix chapter.
Modules \texttt{FastDynamo/Finish/RealSolution.lean}, \texttt{FastDynamo/Operator/Evolution.lean},
\texttt{FastDynamo/Analysis/Realification.lean}.

\begin{lemma}[Exact norm law for the real mode]
  \label{lem:finish:norm-law}
  \lean{FastDynamo.realAxialTrajectory_eq_realPart_complexExp,
  FastDynamo.realAxialInitialField_ne_zero,
  FastDynamo.realAxialTrajectory_ne_zero,
  FastDynamo.norm_realAxialTrajectory,
  FastDynamo.realAxialTrajectory_norm_lower_bound}
  \leanok
  \uses{def:finish:real-solution}
  Let $K\neq0$ and $b\neq0$.  Then $B_{\rm in}\neq0$, $B(t)\neq0$ for every $t\in\R$, and
  \[
    \norm{\re\bigl(\e^{i\omega t}V\bigr)}_2^2=\frac12\norm b_2^2=\norm{\re V}_2^2,
    \qquad\text{hence}\qquad
    \norm{B(t)}_2=\e^{\alpha t}\norm{B_{\rm in}}_2 .
  \]
  If moreover $g\leq\alpha$, then $\e^{gt}\norm{B_{\rm in}}_2\leq\norm{B(t)}_2$ for $t\geq0$.
\end{lemma}

Faithful, and slightly broader than the paper in two respects: the norm identity is proved for
\emph{every} real time, not only $t\geq0$, and the same law is available after transport to the
literal three-torus (\texttt{norm\_literalRealTrajectory\_eq\_exp\_mul\_norm\_zero} in
\texttt{FastDynamo/Analysis/LiteralThreeTorus.lean}).  The auxiliary
\texttt{realAxialTrajectory\_eq\_realPart\_complexExp} identifies the definition of $B(t)$ with
$\re(\e^{t\lambda}V)$.

\begin{proof}
  \leanok
  \uses{lem:real-axial-mode}
  The non-zero axial Fourier index removes the oscillatory cross term when one integrates in
  $x_3$, so no temporal phase can make the real part small.  Apply Lemma
  \ref{lem:real-axial-mode} first with $\varphi=\omega t$ and then with $\varphi=0$: both sides
  equal $\frac12\norm b_2^2$, so the phase factor is norm-preserving and only the real
  exponential $\e^{\alpha t}$ survives.  Non-vanishing at every time follows from the same
  identity.
\end{proof}

\begin{lemma}[The realified mode is a strong solution]
  \label{lem:finish:strong-solution}
  \lean{FastDynamo.DomainOperator.isStrongSolution_eigenTrajectory,
  FastDynamo.productRealPartH2_preserves_divergenceFree,
  FastDynamo.productConjH2_realifiedEigenTrajectoryH2,
  FastDynamo.realifiedEigenTrajectoryH2_divergence_eq_zero,
  FastDynamo.complexify_realAxialTrajectory_eq_realified_embedding,
  FastDynamo.realifiedEigenTrajectoryH2_isPhysicalStrongSolution}
  \leanok
  \uses{def:finish:real-solution, def:finish:physical-operator}
  Let $u$ be divergence-free, let $\eps\in\R$, and let $V\in\Hs^2$ be divergence-free with
  $\mathcal{A}_\eps V=\lambda V$.  Then $W(t)=\re(\e^{t\lambda}V)$ is real (fixed by conjugation)
  and divergence-free at every time, it is the $\Hs^2$ representative of $B(t)$ whenever $V$ is the
  axial lift $\e^{2\pi iKx_3}b$, and it is a strong solution of the physical induction equation
  \[
    \partial_tB=L_{\eps,u}B,\qquad B(0)=B_{\rm in} .
  \]
\end{lemma}

This stands in for the paper's one-line justification in the proof of Theorem~1.2 (``since
the coefficients of the induction equation are real-valued, its solution is
$B^\eps(t)=\e^{\alpha_\eps t}\re(\e^{i\omega_\eps t}V^\eps)$''), which in the paper refers to
the solution given by the semigroup.  Here the complex solution $\e^{t\lambda}V$ is produced directly by scalar
exponentiation (\texttt{DomainOperator.isStrongSolution\_eigenTrajectory}), and no semigroup is
used; its identification with the semigroup orbit is
Lemma~\ref{lem:dynamics:physical-uniqueness}.

\begin{proof}
  \leanok
  \uses{lem:finish:physical-compatibility}
  The curve $t\mapsto\e^{t\lambda}V$ is a strong solution of the auxiliary equation, by
  differentiating the scalar exponential and using the eigenvalue equation.  Taking real parts is a
  fixed continuous real-linear map, so it commutes with the time derivative and with the embeddings
  into $\Lp^2$ and $\Hs^{-2}$; it also preserves the divergence-free constraint.  Finally, for
  divergence-free velocity and divergence-free field the auxiliary and physical actions agree
  (Lemma \ref{lem:finish:physical-compatibility}), which converts the auxiliary equation satisfied
  by the real part into the physical one.
\end{proof}

\begin{proposition}[One unstable reduced mode gives the real growing solution]
  \label{prop:finish:real-strong-mode}
  \lean{FastDynamo.conditionalModeCertificate_of_reducedEigenpair,
  FastDynamo.realStrongModeCertificate_of_reducedEigenpair}
  \leanok
  \uses{def:finish:real-solution, def:main:induction-operator, lem:axial-recovery}
  Let $u$ be divergence-free, $\eps\geq0$, $K\in\Z\setminus\{0\}$, and let $(\lambda,b)$ be an
  eigenpair of the reduced operator $A^{(3)}_{\eps,K}$ with $\re\lambda>0$; let
  $g\leq\re\lambda$.  Then
  the strong-mode certificate of Definition \ref{def:finish:real-solution} holds.  Explicitly, with
  $V=\e^{2\pi iKx_3}b$:
  \begin{itemize}
    \item $V\in\Hs^2(\T^3;\CC^3)$, $\norm V_{\Lp^2(\T^3)}=\norm b_{\Lp^2(\T^2)}>0$, and
      $\mathcal{A}_\eps V=\lambda V$;
    \item $\Div V=0$ and $L_{\eps,u}V=\lambda V$;
    \item $B_{\rm in}=\re V\neq0$ is divergence-free, and $B(t)$ is a real divergence-free strong
      solution of the induction equation with $\norm{B(t)}_2=\e^{t\re\lambda}\norm{B_{\rm in}}_2$
      for all $t\in\R$ and $\e^{gt}\norm{B_{\rm in}}_2\leq\norm{B(t)}_2$ for $t\geq0$.
  \end{itemize}
\end{proposition}

This proposition is the Lean form of the paper's proof of Theorem \ref{thm:main}, minus the
growth-rate clause.  It is the single hinge of the finish layer: the main chapter's
\texttt{MainPhysicalClauses} is obtained from the reduced clauses by applying it (see
\texttt{FastDynamo/Main.lean}).  The remaining clause $\gamma(u,\eps)\geq\gamma_0$ of Theorem
\ref{thm:main} is proved elsewhere and by a different route, from the same reduced eigenpair:
it is Theorem \ref{thm:main-growth-rate}, which needs the $\C_0$-semigroup theory built in
\texttt{FastDynamo/Dynamics/HilleYosida/} and \texttt{FastDynamo/Dynamics/Assembly/} rather
than the analytic-semigroup input of Lemma \ref{lem:semigroup-spectral-growth}.  The norm of the
complex eigentrajectory, $\norm{\e^{t\lambda}V}=\e^{t\re\lambda}\norm V$, is
\texttt{norm\_embedding\_eigenTrajectory}
(\texttt{FastDynamo/Finish/GrowingModeEvolution.lean}).

\begin{proof}
  \leanok
  \uses{lem:positive-eigenfunction-solenoidal, lem:finish:physical-compatibility,
  lem:finish:norm-law, lem:finish:strong-solution, lem:real-axial-mode}
  Since $K\neq0$ and every eigenvector of $A_{\eps,K}$ lies in its domain, the lift
  $V=\e^{2\pi iKx_3}b$ is one-periodic and lies in $\Hs^2$, with
  $\norm V_{\Lp^2(\T^3)}=\norm b_{\Lp^2(\T^2)}$; the axial reduction is an equivalence, so
  $\mathcal{A}_\eps V=\lambda V$.  Because $\re\lambda>0$, Lemma
  \ref{lem:positive-eigenfunction-solenoidal} gives $\Div V=0$, and Lemma
  \ref{lem:finish:physical-compatibility} then gives $L_{\eps,u}V=\lambda V$.  Lemma
  \ref{lem:finish:norm-law} gives the non-vanishing and the exact norm law for
  $B_{\rm in}=\re V$ and $B(t)$, and Lemma \ref{lem:finish:strong-solution} gives the evolution
  equation and the solenoidality of the curve.  Unless $\omega=\im\lambda$ vanishes the datum
  $B_{\rm in}$ is not itself a magnetic eigenmode: it evolves in the real invariant plane spanned
  by $\re V$ and $\im V$.
\end{proof}

\begin{lemma}[Zero mean]
  \label{lem:finish:zero-mean}
  \lean{FastDynamo.HasZeroMean, FastDynamo.productCoeffCLM_zero_zero_eq_integral,
    FastDynamo.hasZeroMean_axialLiftH2, FastDynamo.hasZeroMean_realifiedEigenTrajectoryH2,
    FastDynamo.hasZeroMean_complexify_realAxialTrajectory,
    FastDynamo.hasZeroMean_complexify_realAxialInitialField}
  \leanok
  \uses{def:finish:real-solution}
  Let $K\in\Z\setminus\{0\}$.  For every amplitude $b$ and every $\lambda\in\CC$ the mode
  $V=\e^{2\pi iKx_3}b$, the real datum $B_{\rm in}=\re V$ and the real solution
  $B(t)=\re(\e^{t\lambda}V)$ at every time have zero spatial mean,
  $\int_{\T^3}B(t)\dd x=0$.
\end{lemma}

This is the sentence ``Since $K_\eps\neq0$, every magnetic field constructed here has zero
spatial mean'' of Remark~1.5 of the paper (\texttt{rem:scope-intro}).  The mean is the integral
over the product torus with its probability Haar measure, which is Lebesgue measure on $\T^3$,
and a real field is read in $\CC^3$ through the complexification.  Module
\texttt{FastDynamo/Finish/ZeroMean.lean}.

\begin{proof}
  \leanok
  The mean is the zeroth Fourier coefficient, and the coefficients of $\e^{2\pi iKx_3}b$ are
  supported on the axial frequency $K\neq0$.  Scalar multiples, conjugates and real parts of a
  zero-mean field have zero mean.
\end{proof}

\begin{lemma}[The real part is not an eigenmode]
  \label{lem:finish:real-part-not-eigenmode}
  \lean{FastDynamo.productInductionAction_eq_smul_of_physicalEigenpair,
    FastDynamo.not_isEigenpair_productRealPartH2}
  \leanok
  \uses{def:finish:physical-operator, lem:finish:physical-compatibility,
    def:finish:real-solution}
  Let $u$ be real and divergence-free and let $V\in\Hs^2$ be a divergence-free eigenmode of
  $L_{\eps,u}$ with eigenvalue $\lambda$, $\im\lambda\neq0$.  Then $\re V$ is not an
  eigenfunction of $L_{\eps,u}$, for any eigenvalue.
\end{lemma}

This is Remark~1.4 of the paper (unlabelled): if $\im\lambda_\eps\neq0$, the real part of the
complex magnetic eigenmode $V^\eps$ is not itself a magnetic eigenmode, since it is not an
eigenfunction of $L_{\eps,u}$.  The Lean statement is for every such eigenmode, with
$L_{\eps,u}$ the curl form of Definition~\ref{def:finish:physical-operator}; the divergence-free
hypothesis on $V$ is part of the paper's notion of a magnetic eigenmode.  Module
\texttt{FastDynamo/Finish/RealPartNotEigenmode.lean}.

\begin{proof}
  \leanok
  \uses{lem:finish:physical-compatibility}
  On divergence-free fields $L_{\eps,u}$ is the auxiliary action, which commutes with complex
  conjugation because $u$ is real.  If $\re V=\frac12(V+\bar V)$ were an eigenfunction with
  eigenvalue $\mu$, then $\lambda V+\bar\lambda\bar V=\mu(V+\bar V)$; since
  $\lambda\neq\bar\lambda$, comparing the two relations forces $V=0$, contradicting that $V$
  is an eigenmode.
\end{proof}

\section{Axial Fourier index and magnetic length}

\begin{definition}[Magnetic length of an axial mode]
  \label{def:finish:magnetic-length}
  \lean{FastDynamo.productGradientFromH2,
  FastDynamo.norm_sq_productGradientFromH2_axialLiftH2,
  FastDynamo.complexAxialMagneticLength,
  FastDynamo.phaseAxialGradientL2,
  FastDynamo.phaseAxialGradientL2_norm_sq,
  FastDynamo.complexifyProductR3L2_phaseAxialGradientL2,
  FastDynamo.realPhaseAxialMagneticLength}
  \leanok
  \uses{def:finish:real-solution}
  For $B\in\Hs^2$ write $\nabla B=(\partial_1B,\partial_2B,\partial_3B)$ for the full weak
  gradient, an element of $\Lp^2(\T^3;\CC^3)^3$.  On a fixed axial sector $V=\e^{2\pi iKx_3}b$,
  \[
    \norm{\nabla V}_2^2
    =\norm{\partial_1b}_2^2+\norm{\partial_2b}_2^2+(2\pi K)^2\norm b_2^2 .
  \]
  The $\Hs^1$-based magnetic length of the mode and of its phase-rotated real part are
  \[
    L_B=\frac{\norm V_2}{\norm{\nabla V}_2},
    \qquad
    L_B^{\re,\varphi}
    =\frac{\norm{\re(\e^{i\varphi}V)}_2}{\norm{\re(\e^{i\varphi}\nabla V)}_2},
  \]
  the second one taken at $\varphi=0$ for the real datum $B_{\rm in}$.
\end{definition}

In Lean the real gradient is defined as the vector of the phase-rotated real parts of the three
complex derivative amplitudes (\texttt{phaseAxialGradientL2}), and
\texttt{complexifyProductR3L2\_phaseAxialGradientL2} proves that it is the weak gradient of the
real field: complexified, its components are the weak derivatives of the $\Hs^2$ field
$\re(\e^{i\varphi}V)$, because the weak derivatives commute with conjugation and hence with
the real part.  \texttt{phaseAxialGradientL2\_norm\_sq} proves that
its squared norm is half of the complex one, which is the paper's
$\norm{\nabla B_{\rm in}}_2^2=\frac12\norm{\nabla V}_2^2$ (proof of Corollary~8.3).  Module
\texttt{FastDynamo/Finish/MagneticScale.lean}.

\begin{lemma}[Eigenmode energy estimate]
  \label{lem:finish:magnetic-energy}
  \lean{FastDynamo.ReducedInductionMagneticScale.epsilon_mul_norm_sq_productGradient_le_stretching,
  FastDynamo.ReducedInductionMagneticScale.norm_transverseStretchingCoefficient_pos_of_reducedEigenpair,
  FastDynamo.complexAxialMagneticLength_lower_of_energy,
  FastDynamo.magneticLength_lower_of_energy}
  \leanok
  \uses{def:main:induction-operator, def:finish:magnetic-length}
  Let $u$ be divergence-free, $\eps\geq0$, and let $(\lambda,b)$ be an eigenpair of
  $A^{(3)}_{\eps,K}$ with $\re\lambda\geq0$.  Then, with $M_u$ the $\Lp^\infty$ operator norm of the stretching matrix
  field $b\mapsto(b\cdot\nabla_\perp)u$,
  \[
    \eps\norm{\nabla V}_2^2\leq M_u\norm V_2^2,
    \qquad\text{hence}\qquad
    \sqrt{\eps/M_u}\leq L_B\quad\text{when }M_u>0 .
  \]
  If moreover $\eps>0$, $K\neq0$ and $b\neq0$, then $M_u>0$.
\end{lemma}

This is the energy argument in the proof of Corollary~8.3 of the paper
(\texttt{cor:magnetic-scale}), which the paper does not state as a separate lemma.  One
difference: the paper's constant is
$M_u=\mathrm{ess\,sup}_x\norm{\nabla u(x)}_{\mathrm{op}}$; the Lean constant is the norm of
\texttt{transverseStretchingCoefficient}, the $\Lp^\infty$ operator norm of the stretching matrix
field $x\mapsto Du(x)$, which is what the energy identity uses; by its definition it is the
paper's $M_u$, an identification that no Lean statement records.
Like the paper (``In particular, $M_u>0$'', after the same energy estimate), Lean derives
$M_u>0$ from the existence of the eigenpair, with no property of the constructed velocity.  The hypotheses
$\eps\geq0$, $\re\lambda\geq0$ are weaker than the paper's $\eps>0$, $\re\lambda_\eps>0$.

\begin{proof}
  \leanok
  Take the real part of the eigenvalue identity paired with $V$.  Transverse and axial transport
  are skew, so they drop out; diffusion contributes $-\eps\norm{\nabla V}_2^2$ with the exact
  fixed-sector gradient identity of Definition \ref{def:finish:magnetic-length}; stretching is
  bounded by $M_u\norm V_2^2$.  Since $\re\lambda\geq0$ this gives the energy inequality, and the
  elementary algebraic step $\eps G^2\leq MF^2\Rightarrow\sqrt{\eps/M}\leq F/G$ gives the lower
  bound.  If $\eps>0$ and the gradient is non-zero, the left-hand side is positive, so $M_u>0$.
\end{proof}

\begin{corollary}[Axial Fourier index and $\Hs^1$-based magnetic length of the growing fields]
  \label{cor:magnetic-scale}
  \lean{FastDynamo.complexAxialMagneticLength_le,
  FastDynamo.realPhaseAxialMagneticLength_eq_complex,
  FastDynamo.realPhaseAxialMagneticLength_le,
  FastDynamo.ReducedInductionMagneticScale.complexAxialMagneticLength_lower,
  FastDynamo.ScaleFamily.matching_fourierIndex_sq_bounds,
  FastDynamo.ScaleFamily.matching_fourierIndex_lower,
  FastDynamo.ScaleFamily.matching_fourierIndex_upper,
  FastDynamo.MainPhysicalClausesAt, FastDynamo.mainPhysicalClausesAt_of_reducedEigenpair,
  FastDynamo.main_theorem_magnetic_scale}
  \leanok
  \uses{def:finish:magnetic-length, def:finish:real-solution, def:velocity:scale-family,
    thm:main-spectral, def:main:witnesses}
  Define the $\Hs^1$-based magnetic length of the magnetic eigenmode $V^\eps$ of Theorem
  \ref{thm:main} by $L_B^\eps=\norm{V^\eps}_2/\norm{\nabla V^\eps}_2$.  There are constants
  $c,C>0$, independent of $\eps$, such that
  \[
    c\eps^{-1/2}\leq\abs{K_\eps}\leq C\eps^{-1/2},
    \qquad
    c\sqrt\eps\leq L_B^\eps\leq C\sqrt\eps
    \qquad(0<\eps\leq\eps_0),
  \]
  the range in which Theorem~\ref{thm:main} produces $V^\eps$.  The real datum
  $B_{\rm in}^\eps=\re V^\eps$ has the same magnetic length.
\end{corollary}

The Lean statement is \texttt{main\_theorem\_magnetic\_scale}
(\texttt{FastDynamo/Finish/MagneticScaleMain.lean}): there are $c,C>0$ such that for every
$0<\eps\leq\eps_0$ there is a triple $(K,\lambda,b)$ satisfying the clauses of
Theorem~\ref{thm:main} at the explicit witnesses (\texttt{MainPhysicalClausesAt}, the body of
the sixteen clauses of Definition~\ref{def:main:physical-clauses}) together with the two
two-sided estimates for $\abs K$ and for the magnetic length of the complex mode
$V=\e^{2\pi iKx_3}b$, and such that the real part of $V$ at every phase, in particular the real
datum $B_{\rm in}^\eps$, has the same magnetic length.  The triple is the eigenpair in the axial
sector $K=K_{n(\eps)}$ of the matching cell of the parameter chain.  The general bounds it is
assembled from are also stated on their own:
\begin{itemize}
  \item the index estimate, for every cell $n$ of an arbitrary scale family and every
    $\eps>0$ with $\eps/\ell_n^2\in I$,
    $k_*\sqrt{\delta_-/\eps}\leq K_n\leq k_*\sqrt{\delta_+/\eps}$
    (\texttt{ScaleFamily.matching\_fourierIndex\_lower} and \texttt{\_upper}, with the squared
    form \texttt{matching\_fourierIndex\_sq\_bounds});
  \item the upper magnetic-length bound $L_B\leq1/(2\pi\abs K)$ for every $K\neq0$ and every
    non-zero amplitude, for the complex mode (\texttt{complexAxialMagneticLength\_le}) and at
    every real phase (\texttt{realPhaseAxialMagneticLength\_le});
  \item the lower bound $\sqrt{\eps/M_u}\leq L_B$ of the paper's proof, for every non-zero
    eigenpair of $A^{(3)}_{\eps,K}$ with $\eps>0$, $K\neq0$ and $\re\lambda\geq0$
    (\texttt{ReducedInductionMagneticScale.complexAxialMagneticLength\_lower},
    Lemma~\ref{lem:finish:magnetic-energy});
  \item the identity of magnetic lengths of the real and the complex mode,
    \texttt{realPhaseAxialMagneticLength\_eq\_complex}, for every $K\neq0$, every non-zero
    amplitude and every phase.
\end{itemize}
Deviations: the norm $\norm{\nabla V}_2$ is that of the weak gradient on the product torus and
$\norm{V}_2$ the product-measure norm; the lower bound uses
$M_u=\norm{\texttt{transverseStretchingCoefficient}\ u}$ in place of
$\mathrm{ess\,sup}\norm{\nabla u}_{\rm op}$ (Lemma~\ref{lem:finish:magnetic-energy}); the
constants are $c=\min\{\sqrt{\delta_-},M_u^{-1/2}\}$ and
$C=\max\{\sqrt{\delta_+},(2\pi\sqrt{\delta_-})^{-1}\}$ at $k_*=1$.  The triple $(K,\lambda,b)$
has its own existential in this statement, as in \texttt{MainPhysicalClauses}: that it is the
triple of \texttt{MainSemigroupOrbitClauses} is not stated, and the clause ``its solution'' (the
semigroup orbit and the uniqueness of item (8) of Theorem~\ref{thm:main-strengthened}) is not
part of it.  The remark after the
corollary in the paper (Remark~8.4, unlabelled), which only observes that this magnetic length is
comparable to the classical resistive scale $\sqrt\eps$ of Moffatt and Proctor, has no Lean
counterpart and is not a node.

\begin{proof}
  \leanok
  \uses{lem:real-axial-mode, lem:finish:magnetic-energy, prop:finish:real-strong-mode}
  For the matching index $n=n(\eps)$ one has $\ell_n=\sqrt{\eps/\delta_n}$, $2\pi K_n\ell_n=N$
  and $\delta_n\in[\delta_-,\delta_+]$, so $K_\eps=K_n=k_*\sqrt{\delta_n/\eps}$, which is the
  index estimate.  (The paper writes this arithmetic at $N=2\pi$, where the middle identity
  is $K_n\ell_n=1$ and the index is $K_\eps=\sqrt{\delta_n/\eps}$; Lean keeps
  $k_*$ free, so its constants read $k_*\sqrt{\delta_\mp/\eps}$.)  Since
  $\partial_{x_3}V^\eps=2\pi iK_nV^\eps$, the full gradient dominates
  $2\pi\abs{K_n}\norm{V^\eps}_2$, whence
  $L_B^\eps\leq1/(2\pi K_n)=\ell_n/N=\ell_n/(2\pi k_*)$ -- the paper's
  $L_B^\eps\leq\ell_n/(2\pi)\leq\sqrt\eps/(2\pi\sqrt{\delta_-})$ at $k_*=1$.
  The converse inequality is the energy estimate of Lemma
  \ref{lem:finish:magnetic-energy}.  Finally, applying
  Lemma \ref{lem:real-axial-mode} to $V^\eps$ and to each of its three derivatives gives
  $\norm{B_{\rm in}^\eps}_2^2=\frac12\norm{V^\eps}_2^2$ and
  $\norm{\nabla B_{\rm in}^\eps}_2^2=\frac12\norm{\nabla V^\eps}_2^2$, so the real datum has the
  same magnetic length.  The constants are the minimum and the maximum of the constants of
  the two steps.
\end{proof}
